\tituloI{PLANTEAMIENTO DEL ESTUDIO}

\tituloII{Planteamiento del problema}

\parrafoII{A nivel global, la integración de la inteligencia artificial generativa en el desarrollo de software ha transformado significativamente la forma en que se construyen soluciones algorítmicas, permitiendo automatizar la generación de código y reducir los tiempos de desarrollo. No obstante, investigaciones recientes evidencian que, si bien los modelos de lenguaje de gran escala son capaces de producir soluciones funcionales, estas no siempre presentan un rendimiento computacional óptimo, especialmente en problemas algorítmicos complejos \cite{idrisov2024, niu2024efficiency}. Estudios contemporáneos han demostrado que el código generado por estos modelos puede presentar ineficiencias en términos de tiempo de ejecución y consumo de memoria, lo cual compromete su aplicabilidad en entornos donde la escalabilidad y el uso eficiente de recursos son críticos \cite{elnashar2024performance, coignion2024performance}. Asimismo, la evaluación del código generado por inteligencia artificial continúa siendo un desafío abierto, debido a la falta de métricas estandarizadas y metodologías consistentes que permitan medir su calidad y eficiencia de manera objetiva \cite{ghosh2024benchmarking, tong2024codejudge}.}

\parrafoII{En el contexto nacional, la adopción de herramientas de inteligencia artificial generativa en la educación superior y en la práctica profesional ha crecido de manera acelerada, impulsada por la disponibilidad de modelos de acceso gratuito y su facilidad de uso. Sin embargo, esta adopción no ha estado acompañada de un enfoque riguroso en la evaluación técnica del código generado, lo que ha derivado en una creciente dependencia de soluciones automatizadas sin un adecuado análisis de su rendimiento computacional \cite{indecopi2023}. Esta situación plantea un riesgo significativo, ya que el uso de soluciones tipo “caja negra” puede desplazar el razonamiento algorítmico crítico, limitando la capacidad de los futuros profesionales para identificar ineficiencias y optimizar soluciones en contextos reales \cite{siddiq2023}. En este escenario, la falta de evaluación empírica del rendimiento del código generado por inteligencia artificial representa una debilidad en la formación de ingenieros y en la calidad del software desarrollado.}

\parrafoII{A nivel local, en la ciudad de Cusco y particularmente en el entorno académico de la ingeniería de sistemas, se observa que los estudiantes emplean herramientas de inteligencia artificial generativa para resolver problemas relacionados con estructuras de datos y algoritmos, especialmente en áreas como ordenación y grafos. No obstante, esta práctica se realiza generalmente sin un proceso sistemático de validación del rendimiento computacional de las soluciones generadas, lo que dificulta la identificación de algoritmos subóptimos cuando se enfrentan a grandes volúmenes de datos o condiciones de alta exigencia. Estudios regionales evidencian que la adopción de estas tecnologías en el ámbito educativo no siempre está acompañada de estrategias que promuevan el análisis crítico del código generado, generando una brecha en la formación técnica de los estudiantes \cite{huaman2024}.}

\parrafoII{En este contexto, la unidad de análisis de la presente investigación está constituida por algoritmos de ordenación y grafos generados mediante modelos de inteligencia artificial generativa de acceso gratuito, los cuales son utilizados como soluciones para problemas computacionales específicos. La problemática central radica en la ausencia de evidencia empírica comparativa que permita determinar si dichos algoritmos presentan un rendimiento computacional adecuado en términos de tiempo de ejecución y consumo de memoria frente a implementaciones estándar desarrolladas bajo principios tradicionales de ingeniería de software. Esta falta de análisis sistemático limita la capacidad de validar la eficiencia de las soluciones generadas automáticamente y dificulta su adopción en entornos donde el rendimiento es un factor crítico.}

\parrafoII{La situación descrita se ve agravada por la ausencia de metodologías de benchmarking integradas en el proceso de evaluación del código generado por inteligencia artificial, así como por la limitada incorporación de herramientas de profiling en la formación académica. Como consecuencia, se generan soluciones que, aunque funcionales, pueden presentar un comportamiento ineficiente bajo condiciones de escalabilidad, lo que se traduce en un uso inadecuado de recursos computacionales y una potencial degradación del rendimiento en aplicaciones reales \cite{ghosh2024benchmarking}. Asimismo, la falta de comparación con implementaciones estándar dificulta establecer criterios objetivos que permitan determinar la calidad del código generado por inteligencia artificial.}

\parrafoII{Las implicancias de esta problemática son relevantes tanto en el ámbito académico como profesional. En el contexto educativo, la dependencia de soluciones generadas por inteligencia artificial sin un análisis crítico puede limitar el desarrollo de habilidades fundamentales en el diseño y optimización de algoritmos. En el ámbito profesional, la adopción de código ineficiente puede generar sobrecostos operativos, baja escalabilidad y problemas de rendimiento en sistemas de información. A largo plazo, la ausencia de criterios técnicos para evaluar el rendimiento del código generado por inteligencia artificial puede afectar la competitividad tecnológica y la calidad del software desarrollado.}

\parrafoII{En este sentido, resulta necesario realizar un análisis comparativo del rendimiento computacional de los algoritmos generados por inteligencia artificial generativa frente a implementaciones estándar, mediante el uso de técnicas de benchmarking que permitan medir de manera objetiva el tiempo de ejecución y el consumo de memoria. Esta investigación busca contribuir al desarrollo de criterios técnicos basados en evidencia empírica que permitan evaluar la eficiencia del código generado automáticamente, fortaleciendo el uso responsable de la inteligencia artificial en el desarrollo de software y promoviendo una formación más rigurosa en el ámbito de la ingeniería de sistemas.}
\tituloII{Formulación del Problema}

\tituloIII{Problema General}
\parrafoIII{¿Qué diferencias existen en el rendimiento computacional, en términos de tiempo de ejecución y consumo de memoria, entre los algoritmos generados por inteligencia artificial generativa y las implementaciones estándar?}

\tituloIII{Problemas Específicos}
\begin{listaPE}
    \item ¿Cómo varía el tiempo de ejecución de los algoritmos de ordenación generados por inteligencia artificial frente a implementaciones estándar, considerando diferentes tamaños de entrada?
    \item ¿Qué diferencias existen en el consumo de memoria de los algoritmos de grafos generados por inteligencia artificial en comparación con implementaciones estándar?
    \item ¿Cómo influye el uso de inteligencia artificial generativa en la eficiencia computacional de los algoritmos en comparación con soluciones tradicionales?
\end{listaPE}

\tituloII{Objetivos}

\tituloIII{Objetivo General}
\parrafoIII{Analizar comparativamente el rendimiento computacional de algoritmos de ordenación y grafos generados por inteligencia artificial generativa frente a implementaciones estándar, mediante técnicas de benchmarking.}

\tituloIII{Objetivos Específicos}
\begin{listaOE}
    \item Evaluar el tiempo de ejecución de los algoritmos de ordenación generados por inteligencia artificial en comparación con implementaciones estándar bajo diferentes tamaños de entrada.
    \item Analizar el consumo de memoria de los algoritmos de grafos generados por inteligencia artificial frente a implementaciones estándar.
    \item Comparar el rendimiento computacional global de los algoritmos generados por inteligencia artificial con soluciones tradicionales de referencia.
\end{listaOE}

\tituloII{Justificación de la investigación}

\tituloIII{Justificación Teórica}
\parrafoIII{
La investigación contribuye al análisis empírico del comportamiento de algoritmos generados por modelos de lenguaje, permitiendo contrastar los principios teóricos de la complejidad computacional con resultados experimentales. De esta manera, se aporta evidencia sobre la capacidad real de estas herramientas para producir soluciones eficientes \cite{idrisov2024}.
}

\tituloIII{Justificación Práctica}
\parrafoIII{
Desde una perspectiva aplicada, los resultados permitirán establecer criterios técnicos para la evaluación del código generado por inteligencia artificial, contribuyendo a mejorar la calidad del software desarrollado en entornos académicos y profesionales.
}

\tituloIII{Justificación Metodológica}
\parrafoIII{
La investigación propone un enfoque experimental basado en técnicas de benchmarking y profiling, permitiendo la medición objetiva y replicable de métricas de rendimiento computacional, lo cual puede ser utilizado como referencia en estudios similares.
}

\tituloII{Delimitación de la investigación}

\tituloIII{Delimitación Espacial}
\parrafoIII{
El estudio se desarrollará en un entorno de cómputo controlado utilizando una computadora personal con arquitectura x86\_64 y sistema operativo Linux, así como recursos disponibles en la Universidad Continental, sede Cusco.
}

\tituloIII{Delimitación Temporal}
\parrafoIII{
La ejecución de los experimentos y análisis de resultados se realizará durante el periodo comprendido entre julio y septiembre de 2026.
}

\tituloIII{Delimitación Social}
\parrafoIII{
La investigación se centra en el análisis de algoritmos generados por inteligencia artificial, por lo que no involucra directamente sujetos humanos. No obstante, sus resultados están orientados a la comunidad académica de ingeniería de sistemas.
}

\tituloII{Hipótesis y Variables}

\tituloIII{Hipótesis}
\parrafoIII{
Los algoritmos generados por inteligencia artificial generativa presentan un rendimiento computacional inferior, en términos de tiempo de ejecución y consumo de memoria, en comparación con implementaciones estándar.
}

\tituloIII{Definición de Variables}
\parrafoIII{
La variable principal es el \textbf{Rendimiento Computacional}, evaluada mediante las siguientes dimensiones:
}

\begin{listaVar}
    \item \textbf{Dimensión Temporal:} Tiempo de ejecución del algoritmo, medido en milisegundos (ms).
    \item \textbf{Dimensión Espacial:} Consumo máximo de memoria durante la ejecución, medido en megabytes (MB).
\end{listaVar}